\scanpage{33}%
\mspage{18}\space
qu'il est divisible par $c_4^3-c_6^2$ dans $\mathbb{Z}/(2)[c_4,c_6,\Delta]$ (resp. $\mathbb{Z}/(3)[c_4,c_6,\Delta]$). Cela résulte de la structure de l'anneau des formes modulaires de caractéristique $2$ (resp. $3$).

\uline{Proposition} 7.2. \quad (I) \uline{L'anneau des formes modulaires de caractéristique} $2$ \uline{est} $\mathbb{Z}/(2)[\tikz[baseline=(coeffa.base)]{\node[inner sep=0pt] (coeffa) {$a_1$};\handwritten{\node[inner sep=0pt,anchor=west,font=\scriptsize] (notea) at ([xshift=1.35em,yshift=-1.05em]coeffa.south west) {le diviseur de $a_1$ est lisse sur $\mathbb{F}_2$};\draw[thin] (coeffa.south west) to[out=-86,in=192] ([xshift=-.18em]notea.west);}\pgfresetboundingbox\path[use as bounding box] ([yshift=-1.5em]coeffa.south west) rectangle (coeffa.north east);},\Delta]$ , \uline{et} $c_4=a_1^4$ , $c_6=a_1^6$ .

\hspace*{3em}(II) \uline{L'anneau des formes modulaires de caractéristique} $3$ \uline{est} $\mathbb{Z}/(3)[\tikz[baseline=(coeffb.base)]{\node[inner sep=0pt] (coeffb) {$b_2$};\handwritten{\node[inner sep=0pt,anchor=west,font=\scriptsize] (noteb) at ([xshift=1.35em,yshift=-1.05em]coeffb.south west) {le diviseur de $b_2$ est lisse sur $\mathbb{F}_2$};\draw[thin] (coeffb.south west) to[out=-86,in=192] ([xshift=-.18em]noteb.west);}\pgfresetboundingbox\path[use as bounding box] ([yshift=-1.5em]coeffb.south west) rectangle (coeffb.north east);},\Delta]$ , et $c_4=b_2^2$ , $c_6=-b_2^3$ .

Tout d'abord, il résulte de $(1.6)_1$ et $(1.7)_2$ que $a_1$ (resp. $b_2$) est une forme modulaire de caractéristique $2$ (resp. $3$) .

En caractéristique $2$ , cherchons tout d'abord les formes modulaires relatives aux seules courbes telles que $a_1$ soit inversible. Ces courbes se mettent sous forme (cf. (4.\tikz[baseline=(un.base)]{\node[inner sep=0pt] (un) {1};\handwritten{\draw[thin,gray,rotate around={-12:(un.center)}] (un.center) ellipse (0.36em and 1.12em);}}))
\[
 y^2+a_1xy=x^3+a_2x^2+a_6 \qquad \pi=\frac{dx}{a_1x}
\]
et, compte tenu de (4.3), on voit que ces formes s'écrivent
\[
\begin{aligned}
 &a_1^{-n}\,P(a_1,a_6) \qquad \text{pour }P\text{ isobare, ou encore}\\[.65em]
 &a_1^{-m}\,Q(a_1,\Delta)\quad\handwritten{,}
\end{aligned}
\]
et cela garde un sens pour $a_1=0$ , $\Delta\ne0$ si et seulement si on avait en fait affaire à un polyn\tikz[baseline=(oc.base)]{\node[inner sep=0pt] (oc) {ô};\handwritten{\draw[thin,gray,rotate around={-14:(oc.center)}] (oc.center) ellipse (0.56em and 1.1em);}}me en $a_1$ et $\Delta$ .

En caractéristique $3$ , raisonnant de même d'abord sur les courbes telles que $c_4$ soit inversible, et remarquant que $d_6=0$ équivaut alors à $a_4=0$ dans (3.1), on voit que toute forme modu%
% archive sheet number: 32
